% INTRODUCCIÓN

\section{Introducción}

Descripción y contextualización del Trabajo Fin de Máster. Si crees que merece la pena, cuenta las motivaciones que te han llevado a realizar este trabajo. Como ejemplo de cita en formato APA: \cite{goodfellow2014generative}.

Describe la estructura general del documento.

El estilo del párrafo tiene que ser encuadrado.

En caso de terminar con un apartado (nivel uno de título) aplicar un salto de página para comenzar siempre el siguiente apartado en una página nueva. No tener en cuenta si es par o impar para el comienzo de un apartado.

Tanto las figuras como las tablas tienen que estar indicadas en el texto, referenciadas según la numeración que se tiene. La descripción debe ser clara y explicar lo que se quiere representar con independencia de la referencia del texto.

Tanto las tablas como las figuras tendrán un estilo de párrafo centrado. La descripción se realizará desde “Insertar título”.

\begin{center}
\begin{longtblr}[entry=none,label=none]{colspec={l l l}}
\textbf{operación} & \textbf{símbolo} & \textbf{ejemplo} \\
suma & + & 4 + 4 = 8 \\
resta & - & 4 - 2 = 2 \\
multiplicación & * & 4 * 4 = 16 \\
\end{longtblr}
\end{center}
\captionof{table}{Operaciones matemáticas utilizadas en el estudio realizado. Elaboración propia.}
\label{tab:operaciones}
